\documentclass[../../main.tex]{subfiles}

\begin{document}

\chapter{Kriptografi Kurva Eliptik (ECC)}
\label{chap:ecc}

\begin{learningoutcome}
Setelah mempelajari bab ini, mahasiswa diharapkan mampu:
\begin{itemize}
    \item Menjelaskan konsep dasar aljabar abstrak (grup, medan, dan Medan Galois) yang mendasari ECC (Sub-CPMK 1.3).
    \item Menganalisis properti geometris dan analitik dari kurva eliptik, termasuk operasi penjumlahan dan penggandaan titik (Sub-CPMK 1.3).
    \item Menguraikan mekanisme kerja \textit{Elliptic Curve Discrete Logarithm Problem} (ECDLP) sebagai basis keamanan ECC (Sub-CPMK 1.3).
    \item Menerapkan prosedur pertukaran kunci ECDH dan enkripsi ECEG (Sub-CPMK 1.3).
    \item Mengevaluasi efisiensi ECC dibandingkan dengan RSA dari sisi panjang kunci dan kebutuhan sumber daya komputasi (Sub-CPMK 1.3).
    \item Menjelaskan metode \textit{encoding} pesan menjadi titik pada kurva eliptik (Sub-CPMK 1.3).
    \item Menerapkan aritmetika medan berhingga, baik $GF(p)$ maupun $GF(2^m)$, untuk menghitung pada kurva eliptik (Sub-CPMK 1.3).
    \item Memeriksa orde grup dan orde titik basis untuk menilai kelayakan sebuah kurva eliptik (Sub-CPMK 1.3).
\end{itemize}
\end{learningoutcome}

\subfile{section-01}
\subfile{section-02}
\subfile{section-03}
\subfile{section-04}
\section{Contoh Terhitung}
\label{sec:contoh-ecc}

Contoh-contoh pada bagian ini disusun mengikuti alur bab, dari aritmetika
medan sampai protokol yang lengkap. Contoh pertama bekerja di dalam
$GF(2^8)$, tempat seluruh aturan medan berhingga terlihat paling jelas.
Contoh kedua dan ketiga menjalankan penjumlahan serta penggandaan titik pada
kurva, masing-masing di atas medan berhingga dan di atas bilangan riil.
Tiga contoh terakhir merangkai seluruh perkakas itu menjadi protokol
pertukaran kunci, enkripsi, dan pengkodean pesan. Setiap angka di bawah ini
telah diverifikasi dengan program pada bagian praktikum di akhir bab ini.

\begin{example}[Aritmetika polinomial di dalam medan $GF(2^8)$]
\label{ex:gf2m}
Setiap \textit{byte} dipandang sebagai polinom berderajat paling tinggi tujuh,
dengan setiap bit menjadi koefisien. Dengan kesepakatan itu,
\[
\code{0D} = x^3 + x^2 + 1,
\qquad
\code{06} = x^2 + x .
\]

\textbf{(a) Penjumlahan.} Penjumlahan di dalam $GF(2)$ dilakukan modulo dua,
sehingga $1 + 1 = 0$; koefisien yang sama pada kedua polinom saling
menghapus. Karena itu penjumlahan dua polinom tidak lain adalah operasi XOR
bit demi bit. Untuk polinom di atas,
\[
(x^3 + x^2 + 1) \oplus (x^2 + x) = x^3 + (x^2 \oplus x^2) + x + 1
= x^3 + x + 1 = \code{0B} .
\]
Suku $x^2$ lenyap karena muncul pada kedua polinom. Dengan cara yang sama,
\[
\code{57} \oplus \code{83} = \code{D4} ,
\]
sebab dalam biner $\code{57} = 01010111$, $\code{83} = 10000011$, dan
XOR keduanya adalah $11010100 = \code{D4}$.

\textbf{(b) Perkalian.} Perkalian dilakukan dalam dua tahap: kalikan dahulu
seperti polinom biasa, lalu reduksi hasilnya dengan polinom reduksi
$x^8 + x^4 + x^3 + x + 1$ (yaitu $\code{11B}$, polinom yang dipakai AES).
Tulis kedua faktor dalam bentuk polinom,
\[
\code{57} = x^6 + x^4 + x^2 + x + 1,
\qquad
\code{83} = x^7 + x + 1 .
\]
Karena $\code{83}$ memiliki bit pada posisi $7$, $1$, dan $0$, perkaliannya
adalah XOR dari tiga salinan $\code{57}$ yang digeser sejauh posisi bit
tersebut. Tabel~\ref{tab:gf2m-kali} memperlihatkan langkah itu.

\begin{table}[htbp]
    \centering
    \caption{Perkalian dua polinom di dalam $GF(2^8)$: menghasilkan polinom
    mentah berderajat $13$, lalu direduksi menjadi berderajat $7$}
    \label{tab:gf2m-kali}
    \begin{tabular}{@{}l l l@{}}
    \toprule
    \textbf{Langkah} & \textbf{Biner} & \textbf{Bentuk polinom} \\
    \midrule
    $\code{57} \ll 0$        & $01010111$          & $x^6 + x^4 + x^2 + x + 1$ \\
    \addlinespace
    $\code{57} \ll 1$        & $10101110$          & $x^7 + x^5 + x^3 + x^2 + x$ \\
    \addlinespace
    $\code{57} \ll 7$        & $10101110000000$    & $x^{13} + x^{11} + x^9 + x^8 + x^7$ \\
    \midrule
    XOR ketiganya            & $10101101111001$    & $x^{13} + x^{11} + x^9 + x^8 + x^6
                             + x^5 + x^4 + x^3 + 1$ \\
    \midrule
    $\oplus\,(\code{11B} \ll 5)$ & $0000100000011001$ & $x^{11} + x^4 + x^3 + 1$ \\
    \addlinespace
    $\oplus\,(\code{11B} \ll 3)$ & $11000001$         & $x^7 + x^6 + 1 = \code{C1}$ \\
    \bottomrule
    \end{tabular}
\end{table}

Baris keempat Tabel~\ref{tab:gf2m-kali} adalah hasil kali mentah, yang
berderajat $13$ sehingga masih di luar medan. Dua baris terakhir
memperlihatkannya direduksi: selama derajatnya paling sedikit delapan,
polinom reduksi digeser ke posisi bit terdepan lalu di-XOR-kan. Hasil
akhirnya adalah $x^7 + x^6 + 1$, yaitu $\code{C1}$. Jadi
\[
\code{57} \cdot \code{83} = \code{C1} .
\]

\textbf{(c) Pemeriksaan pada medan yang lebih kecil.} Dengan cara yang sama,
di dalam $GF(2^4)$ yang memakai polinom reduksi $x^4 + x + 1$,
\[
(x^3 + x^2 + 1)(x^2 + x) = x^5 + x^3 + x^2 + x ,
\]
sebab koefisien $x^4$ dan $x^3$ yang muncul dua kali saling menghapus. Hasil
kali itu berderajat lima, sehingga direduksi satu langkah:
\[
(x^5 + x^3 + x^2 + x) \oplus \bigl(x \cdot (x^4 + x + 1)\bigr)
= (x^5 + x^3 + x^2 + x) \oplus (x^5 + x^2 + x) = x^3 .
\]
Sisanya berderajat tiga, jadi hasil akhirnya adalah $x^3 = \code{1000}$.

\par\vspace{2pt}{\footnotesize\itshape Catatan: dua angka pada berkas sumber
perlu dikoreksi. Pertama, berkas sumber menuliskan hasil kali mentah di atas
sebagai $\code{10110}$, padahal bentuk binernya adalah $\code{101110}$; dua
bit terdepan $\code{10}$ hilang. Kedua, berkas sumber menyebut $x^2 + 1$
tidak dapat direduksi, padahal di dalam $GF(2)$ polinom itu sama dengan
$(x+1)^2$, sebab $(x+1)^2 = x^2 + 2x + 1 = x^2 + 1$. Karena itu berkas
sumber juga salah ketika memfaktorkan $x^5 + x^2 + x + 1$ menjadi
$(x^3 + x^2 + 1)(x^2 + 1)$; hasil kali kedua faktor itu adalah
$x^5 + x^4 + x^3 + 1$, bukan $x^5 + x^2 + x + 1$. Faktorisasi yang benar
adalah $x^5 + x^2 + x + 1 = (x^2 + 1)(x^3 + x + 1)$. Kedua koreksi itu
diperiksa oleh program pada bagian praktikum.\par}
\end{example}

\begin{example}[Penjumlahan titik pada kurva di atas $GF(11)$]
\label{ex:gf11}
Ambil kurva
\[
y^2 \equiv x^3 + x + 6 \pmod{11},
\]
yang ordonya $13$ dan telah dienumerasi pada Subbagian~\ref{subsec:enumerasi-gf-p}.
Dua di antara titiknya adalah $P(2, 4)$ dan $Q(5, 9)$; keduanya memenuhi
kurva, sebab $4^2 = 16 \equiv 5$ dan $2^3 + 2 + 6 = 16 \equiv 5$ untuk $P$,
serta $9^2 = 81 \equiv 4$ dan $5^3 + 5 + 6 = 136 \equiv 4$ untuk $Q$.

\textbf{(a) Penjumlahan $P + Q$.} Dengan $a = 1$ dan $p = 11$, gradiennya
menurut Persamaan~\eqref{eq:gradien-mod-p} adalah
\[
m \equiv (9 - 4)(5 - 2)^{-1} \equiv 5 \cdot 3^{-1} \pmod{11} .
\]
Karena $3 \cdot 4 = 12 \equiv 1$, invers dari $3$ adalah $4$, sehingga
$m \equiv 5 \cdot 4 = 20 \equiv 9$. Selanjutnya
\[
x_r \equiv m^2 - x_p - x_q \equiv 81 - 2 - 5 = 74 \equiv 8 \pmod{11},
\]
\[
y_r \equiv m(x_p - x_r) - y_p \equiv 9(2 - 8) - 4 = -58 \equiv 8 \pmod{11},
\]
sebab $-58 + 6 \cdot 11 = 8$. Jadi $P + Q = (8, 8)$, dan pemeriksaannya:
$8^2 = 64 \equiv 9$ sedangkan $8^3 + 8 + 6 = 526 = 47 \cdot 11 + 9$. Cocok.

\textbf{(b) Penggandaan $2P$.} Rumusnya sama, hanya gradiennya diganti
menurut Persamaan~\eqref{eq:gradien-penggandaan}:
\[
m \equiv (3 \cdot 2^2 + 1)(2 \cdot 4)^{-1} \equiv 13 \cdot 8^{-1}
\pmod{11} .
\]
Karena $8 \cdot 7 = 56 \equiv 1$, invers dari $8$ adalah $7$, sehingga
$m \equiv 13 \cdot 7 = 91 \equiv 3$. Koordinat hasilnya
\[
x_r \equiv m^2 - 2x_p \equiv 9 - 4 = 5, \qquad
y_r \equiv m(x_p - x_r) - y_p \equiv 3(2 - 5) - 4 = -13 \equiv 9 ,
\]
jadi $2P = (5, 9)$.

Perhatikan hasil yang menarik: $2P$ ternyata sama dengan $Q$. Artinya $Q$
bukan titik yang bebas, melainkan kelipatan dari $P$, dan karena itu
$P + Q = 3P = (8, 8)$. Kebetulan semacam ini lazim pada kurva kecil dan
menjadi alasan mengapa pada kurva yang dipakai sungguhan orderya harus
diperiksa lebih dahulu. Bila titik basis berorde kecil, hubungan seperti ini
memberi petunjuk kepada penyerang. Tabel pelelaran $kP$ untuk
$k = 1, 2, \dots, 13$ pada kurva ini, yang berakhir di $13P = \mathcal{O}$,
diperlihatkan pada Tabel~\ref{tab:pelelaran-gf11}.
\end{example}

\begin{example}[Penjumlahan dan penggandaan titik di atas bilangan riil]
\label{ex:real}
Kurva di atas bilangan riil dipakai untuk memahami geometrinya, bukan untuk
kriptografi. Dua kurva berikut memperlihatkan kedua operasi itu.

\textbf{(a) Kurva $y^2 = x^3 - 3x + 3$.} Di sini $a = -3$ dan $b = 3$,
sehingga $4(-3)^3 + 27(3)^2 = -108 + 243 = 135 \neq 0$ menurut
Persamaan~\eqref{eq:diskriminan}: kurva itu licin. Titik $P(1, 1)$ memenuhi
kurva karena $1^2 = 1 - 3 + 3 = 1$, dan titik $Q(-2, 1)$ juga karena
$1^2 = -8 + 6 + 3 = 1$.

Untuk $P + Q$ dengan $P \neq Q$,
\[
m = \frac{1 - 1}{-2 - 1} = \frac{0}{-3} = 0 ,
\]
sehingga garis yang melalui $P$ dan $Q$ mendatar. Akibatnya
\[
x_r = 0 - 1 - (-2) = 1, \qquad y_r = 0 \cdot (1 - 1) - 1 = -1 ,
\]
dan $P + Q = (1, -1)$. Titik itu benar pada kurva sebab
$(-1)^2 = 1 = 1 - 3 + 3$.

Untuk $2P$,
\[
m = \frac{3(1)^2 - 3}{2(1)} = \frac{0}{2} = 0 ,
\]
sehingga garis singgungnya pun mendatar. Maka $x_r = 0 - 2 = -2$ dan
$y_r = 0 \cdot (1 + 2) - 1 = -1$, jadi $2P = (-2, -1)$. Periksa:
$(-1)^2 = (-2)^3 - 3(-2) + 3 = -8 + 6 + 3 = 1$. Cocok.

\textbf{(b) Kurva $y^2 = x^3 + 2x + 4$.} Di sini $a = 2$ dan $b = 4$.
Titik $P(2, 4)$ memenuhi kurva karena $4^2 = 8 + 4 + 4 = 16$, dan
$Q(0, 2)$ juga karena $2^2 = 0 + 0 + 4 = 4$. Gradien garis $PQ$ adalah
\[
m = \frac{2 - 4}{0 - 2} = \frac{-2}{-2} = 1 ,
\]
sehingga $x_r = 1 - 2 - 0 = -1$ dan $y_r = 1 \cdot (2 + 1) - 4 = -1$.
Jadi $P + Q = (-1, -1)$, dan titik itu benar pada kurva sebab
$(-1)^2 = 1 = -1 - 2 + 4$.

Penggandaannya menghasilkan pecahan:
\[
m = \frac{3(2)^2 + 2}{2(4)} = \frac{14}{8} = \frac{7}{4},
\]
\[
x_r = \left(\frac{7}{4}\right)^2 - 2(2) = \frac{49}{16} - 4 = -\frac{15}{16},
\]
\[
y_r = \frac{7}{4}\left(2 + \frac{15}{16}\right) - 4
    = \frac{7 \cdot 47}{64} - 4 = \frac{329 - 256}{64} = \frac{73}{64} .
\]
Jadi $2P = \left(-\frac{15}{16}, \frac{73}{64}\right)$, atau dalam pecahan
desimal $(-0{,}9375;\ 1{,}140625)$. Bandingkan dengan kurva di atas
$GF(11)$: di sana koordinat hasilnya tetap bilangan bulat, sedangkan di sini
koordinatnya langsung menjadi pecahan. Itulah sebabnya kriptografi memakai
kurva di atas medan berhingga.
\end{example}

\begin{example}[Pertukaran kunci ECDH]
\label{ex:ecdh}
\textbf{(a) Pada kurva di atas $GF(5)$.} Ambil kurva
\[
y^2 \equiv x^3 + 2x + 1 \pmod{5},
\]
yang memiliki $6$ titik berhingga ditambah $\mathcal{O}$, sehingga orderya
$7$. Titik basis yang disepakati adalah $B(0, 1)$, yang memenuhi kurva sebab
$1^2 = 0 + 0 + 1$. Alice memilih kunci privat $a = 2$ dan Bob $b = 3$.
Kunci publik masing-masing adalah
\[
P_A = 2B = (1, 3),
\qquad
P_B = 3B = (3, 3).
\]
Setiap pihak menghitung kunci bersama dari kunci publik pihak lain:
\[
K = a \cdot P_B = 2(3, 3) = (0, 4),
\qquad
K = b \cdot P_A = 3(1, 3) = (0, 4) .
\]
Keduanya memperoleh titik $(0, 4)$ --- perhatikan bahwa titik itu sama dengan
$-B$, sebab $B$ berorde $7$ sehingga $6B = -B$.

\textbf{(b) Pada kurva di atas $GF(23)$.} Kurva $y^2 \equiv x^3 + x + 1
\pmod{23}$ dengan titik basis $B(3, 10)$ memberi angka yang lebih besar.
Dengan $a = 2$ dan $b = 3$,
\[
P_A = 2B = (7, 12),
\qquad
P_B = 3B = (19, 5),
\]
\[
K = 2P_B = (12, 4),
\qquad
K = 3P_A = (12, 4) .
\]
Penyerang yang menyadap pertukaran itu hanya melihat kurva, titik basis
$B(3,10)$, dan kedua kunci publik $(7,12)$ serta $(19,5)$. Untuk menemukan
$K$ ia harus memulihkan $a$ atau $b$ dari $P_A = aB$, yaitu memecahkan
persoalan logaritma diskrit kurva eliptik pada Persamaan~\eqref{eq:ecdlp}.
Pada kurva mainan ini penyerangan itu sepele --- program pada bagian
praktikum melakukannya dengan penelusuran langkah bayi-langkah raksasa ---
tetapi pada kurva 256 bit usaha yang sama menuntut sekitar $2^{128}$ langkah.
\end{example}

\begin{example}[Enkripsi dan dekripsi ECEG]
\label{ex:eceg}
Enkripsi kurva eliptik memakai kurva yang sama dengan contoh sebelumnya, yaitu
$y^2 \equiv x^3 + x + 1 \pmod{23}$ dengan titik basis $B(3, 10)$. Kunci
privat penerima adalah $b = 3$, sehingga kunci publiknya
\[
P_B = 3B = (19, 5) .
\]
Pesan yang hendak dikirim telah dikodekan lebih dahulu menjadi sebuah titik
pada kurva (caranya diperlihatkan pada contoh berikutnya); di sini misalkan
$P_M = (7, 12)$.

Pengirim memilih kunci sesaat $k = 5$ --- bilangan acak yang baru untuk
setiap pesan --- lalu menghitung dua titik menurut
Persamaan~\eqref{eq:eceg-cipher}:
\[
C_1 = kB = 5B = (9, 16),
\]
\[
C_2 = P_M + k P_B = (7, 12) + 5(19, 5) = (7, 12) + (1, 16) = (18, 3) .
\]
Jadi cipherteksnya adalah pasangan titik $\bigl((9, 16),\ (18, 3)\bigr)$.

Penerima memegang $b = 3$ dan memulihkan pesannya dengan mengurangi titik
kedua dengan hasil kali kunci privatnya dan titik pertama:
\[
P_M = C_2 - b\,C_1 = (18, 3) - 3(9, 16) .
\]
Karena $C_1 = 5B$, maka $3C_1 = 15B$, dan $15B = (1, 16)$ menurut kenyataan
bahwa $B$ berorde $28$. Lawan titik itu adalah $(1, 7)$, sebab
$-16 \equiv 7 \pmod{23}$, sehingga
\[
P_M = (18, 3) + (1, 7) = (7, 12) .
\]
Plainteksnya pulih tepat seperti semula.

Perhatikan mengapa langkah itu selalu berhasil. Penerima tidak perlu mengetahui
$k$; ia hanya perlu bahwa $kP_B = kbB = bkB = bC_1$, sehingga suku acak pada
$C_2$ terhapus oleh suku yang ia hitung sendiri --- persis seperti penurunan
pada Persamaan~\eqref{eq:eceg-dekripsi}. Sekalipun penyerang menyadap kedua
titik cipherteks, ia tidak dapat melakukan hal yang sama karena tidak memegang
$b$.

Satu hal terakhir: karena $k$ dipilih baru untuk setiap pesan, titik
$C_1 = kB$ juga berubah setiap kali. Pesan $P_M$ yang sama karena itu
menghasilkan cipherteks yang berbeda pada setiap pengiriman. Sifat itu,
yang disebut enkripsi probabilistik, membuat penyerang tidak dapat mengenali
pesan yang berulang --- keunggulan yang tidak dimiliki buku kode maupun
pemetaan karakter langsung pada Subbagian~\ref{subsec:pemetaan-langsung}.
\end{example}

\begin{example}[Pengkodean pesan dengan metode Koblitz]
\label{ex:koblitz}
Agar pesan dapat dienkripsi dengan ECEG, rangkaian bitnya harus diubah
terlebih dahulu menjadi titik pada kurva. Metode Koblitz melakukan itu dengan
memanfaatkan peluang: ruas kanan kurva hanya punya peluang sekitar setengah
untuk menjadi kuadrat modulo $p$ pada nilai $x$ yang dipilih acak.

Ambil kurva
\[
y^2 \equiv x^3 - x + 188 \pmod{751},
\]
yang memiliki $N = 727$ titik. Menurut kesepakatan pengkodean pada
Subbagian~\ref{subsec:koblitz}, huruf \texttt{A} sampai \texttt{Z} bernilai
$10$ sampai $35$, sehingga huruf \texttt{B} bernilai $m = 11$. Ambil
parameter basis $k = 20$, sehingga nilai $x$ pertama yang dicoba adalah
\[
x = mk + 1 = (11)(20) + 1 = 221
\]
menurut Persamaan~\eqref{eq:koblitz-x}. Sulihkan ke dalam kurva:
\[
221^3 - 221 + 188 = 10\,793\,861 - 33 = 10\,793\,828
\equiv 456 \pmod{751} .
\]
Apakah $456$ kuadrat modulo $751$? Pemeriksaan pada program menunjukkan tidak,
sehingga $x = 221$ tidak memberi titik. Pencarian berlanjut seperti pada
Tabel~\ref{tab:koblitz-b}: $x = 222$ memberi $446$ dan $x = 223$ memberi
$266$, keduanya juga bukan kuadrat. Baru pada percobaan keempat berhasil.
Untuk $x = 224$ diperoleh
\[
224^3 - 224 + 188 = 11\,239\,424 - 36 = 11\,239\,388
\equiv 673 \pmod{751} ,
\]
dan $673$ adalah kuadrat modulo $751$: dua akarnya adalah $y = 248$ dan
$y = 503$, sebab $248^2 = 61\,504 = 81 \cdot 751 + 673$ dan
$503^2 = 253\,009 = 336 \cdot 751 + 673$. Jadi huruf \texttt{B} dikodekan
menjadi titik $(224, 248)$ pada kurva tersebut.

Penerima memulihkan pesannya dengan Persamaan~\eqref{eq:koblitz-dekode},
\[
m = \left\lfloor \frac{224 - 1}{20} \right\rfloor
  = \left\lfloor \frac{223}{20} \right\rfloor
  = \lfloor 11{,}15 \rfloor = 11 ,
\]
yang memang nilai untuk huruf \texttt{B}.

Empat percobaan pada contoh ini sejalan dengan dugaan peluangnya. Setiap
percobaan berpeluang sekitar setengah untuk berhasil, sehingga rata-rata
dibutuhkan dua percobaan, dan peluang menuntut lebih dari sepuluh percobaan
lebih kecil dari $10^{-3}$. Agar pemetaan itu tetap satu-satu, kurva harus
memiliki paling sedikit $36k$ titik, sebab setiap nilai $m$ memakai wilayah
$x$ selebar $k$ dan ada $36$ nilai $m$ yang mungkin.
\end{example}


\begin{obeactivity}
\textbf{Aktivitas 14.1: Simulasi Penjumlahan Titik}
1. Gunakan kurva $y^2 = x^3 + 2x + 4$ di $\mathbb{R}$.
2. Diberikan titik $P(2, 4)$ dan $Q(0, 2)$, hitunglah koordinat titik $R = P + Q$.
3. Lakukan penggandaan titik untuk menghitung $2P$ pada kurva yang sama.
4. Bandingkan hasil perhitungan manual Anda dengan kalkulator ECC online.
\end{obeactivity}


\begin{obereflection}
\textbf{Latihan 14.1: Aritmetika Medan dan Struktur Aljabar}
1. Mengapa titik tak hingga $\mathcal{O}$ diperlukan dalam definisi grup kurva eliptik?
2. Jelaskan mengapa ECDLP dianggap lebih sulit dipecahkan daripada masalah pemfaktoran bilangan bulat pada RSA.
3. Apa dampak penggunaan Medan Galois $\text{GF}(p)$ dibandingkan bilangan riil dalam implementasi praktis ECC?
4. Perhatikan himpunan bilangan bulat $\Z$ dengan operasi perkalian. Menurut
   Tabel~\ref{tab:contoh-grup}, susunan itu bukan grup. Sebutkan aksioma mana
   yang gagal dan berikan satu contoh penyangkalnya.
5. Tentukan invers perkalian dari $3$, $8$, dan $20$ modulo $11$. Tuliskan pula
   seluruh pasangan $(a, a^{-1})$ untuk $a = 1, 2, \dots, 10$ modulo $11$, lalu
   jelaskan mengapa setiap barisnya pasti ada
   (Subbagian~\ref{subsec:medan-berhingga}).
6. Selesaikan $x^2 \equiv 5 \pmod{11}$ dan $x^2 \equiv 7 \pmod{11}$. Persamaan
   mana yang tidak punya penyelesaian, dan bagaimana kenyataan itu dipakai
   oleh metode Koblitz?
7. Hitung hasil ketiga operasi berikut di dalam $GF(2^8)$ dengan polinom
   reduksi $\code{11B}$: $\code{0D} + \code{06}$, $\code{57} + \code{83}$, dan
   $\code{57} \cdot \code{83}$. Periksa jawaban Anda terhadap
   Tabel~\ref{tab:gf2m-kali}.
8. Hitung $\code{57} \cdot \code{02}$ di dalam medan yang sama. Mengapa
   hasilnya cukup diperoleh dengan menggeser satu bit, tanpa perlu reduksi?
9. Periksa apakah $x^4 + x + 1$ dan $x^2 + 1$ dapat direduksi di dalam $GF(2)$.
   Mengapa hanya polinom yang tidak dapat direduksi yang boleh dipakai
   sebagai polinom reduksi?
\end{obereflection}

\begin{obereflection}
\textbf{Latihan 14.2: Geometri dan Aritmetika Titik}
1. Pada kurva $y^2 \equiv x^3 + x + 6 \pmod{11}$, hitunglah $(2,4) + (8,8)$ dan
   $2(8,8)$. Periksa hasil Anda dengan Tabel~\ref{tab:pelelaran-gf11}.
2. Tunjukkan bahwa titik $(10, 9)$ \emph{tidak} terletak pada kurva
   $y^2 \equiv x^3 + x + 1 \pmod{23}$. Mana yang gagal, ruas kiri atau ruas
   kanan, dan berapa selisihnya?
3. Untuk titik basis $B(3,10)$ pada kurva mod 23, hitunglah $2B$, $3B$, $6B$,
   dan $14B$. Titik manakah yang berordinat nol, dan berapa ordenya?
4. Lanjutkan latihan sebelumnya: hitung $2 \cdot (14B)$, yaitu $28B$. Mengapa
   hasilnya $\mathcal{O}$, dan apa kaitannya dengan orde titik $(4,0)$?
5. Diketahui orde grup kurva mod 23 adalah $28$. Tentukan orde titik
   $(0,1)$ dan $(13,16)$ dengan menghitung kelipatan $kP$ sampai diperoleh
   $\mathcal{O}$. Apakah kedua titik itu dapat dipakai sebagai titik basis?
6. Pada kurva riil $y^2 = x^3 + 2x + 4$, hitunglah $2Q$ untuk $Q(0,2)$. Nyatakan
   hasilnya sebagai pecahan, lalu jelaskan mengapa koordinatnya tidak lagi
   berupa bilangan bulat (Subbagian~\ref{subsec:mengapa-medan}).
7. Tentukan gradien garis singgung kurva $y^2 = x^3 - 4x$ di titik $(-2, 0)$.
   Mengapa nilainya tidak terdefinisi, dan apa artinya bagi hasil $P + P$
   (Subbagian~\ref{subsec:penggandaan-titik})?
8. Periksa syarat ketak-singularan pada
   Persamaan~\eqref{eq:diskriminan} bagi pasangan $(a,b) = (-4, 0)$,
   $(0, 0)$, dan $(-3, 3)$. Kurva mana yang licin?
\end{obereflection}

\begin{obereflection}
\textbf{Latihan 14.3: Protokol, Pemetaan Pesan, dan Efisiensi}
1. Jalankan pertukaran kunci ECDH pada kurva $y^2 \equiv x^3 + x + 1 \pmod{23}$
   dengan titik basis $B(3,10)$, kunci privat Alice $a = 4$, dan kunci privat
   Bob $b = 9$. Tentukan kedua kunci publik dan kunci bersamanya, lalu
   periksa hasilnya dari kedua arah menurut Persamaan~\eqref{eq:ecdh-kunci}.
2. Ulangi pertukaran itu dengan $a = 5$ dan $b = 7$. Mengapa kunci bersamanya
   kebetulan sama dengan salah satu kunci publik? Jelaskan dengan menghitung
   $ab$ modulo orde grup.
3. Lakukan enkripsi ECEG pada kurva yang sama dengan kunci publik penerima
   $P_B(19,5)$, plainteks $P_M(7,12)$, dan kunci sesaat $k = 6$. Tuliskan
   cipherteksnya, lalu pulihkan plainteksnya mengikuti
   Persamaan~\eqref{eq:eceg-dekripsi}.
4. Seandainya penyerang menyadap cipherteks pada latihan sebelumnya. Informasi
   apa yang harus ia temukan lebih dahulu agar dapat memulihkan $P_M$, dan
   berapa biaya penyerangan itu pada kurva 256 bit?
5. Kodekan huruf \texttt{K} (bernilai $m = 20$) pada kurva
   $y^2 \equiv x^3 - x + 188 \pmod{751}$ dengan parameter basis $k = 20$.
   Berapa nilai $x$ yang diperoleh, dan berapa percobaan yang diperlukan?
   Periksa hasilnya dengan menjalankannya pada program.
6. Ulangi untuk huruf \texttt{R} (bernilai $m = 27$). Berapa percobaan yang
   diperlukan kali ini? Bandingkan dengan hasil pada latihan sebelumnya, lalu
   jelaskan mengapa jumlah percobaan berbeda walaupun nilai $m$-nya hanya
   berselisih tujuh.
7. Pulihkan nilai $m$ dari absis titik yang diperoleh pada kedua latihan di
   atas dengan Persamaan~\eqref{eq:koblitz-dekode}. Apakah hasilnya kembali
   menjadi $20$ dan $27$?
8. Untuk menyalurkan kunci sesi AES-128, RSA memerlukan modulus $3072$ bit
   sedangkan ECC cukup $256$ bit. Hitung perbandingan keduanya, lalu jelaskan
   dengan Tabel~\ref{tab:ecdlp-vs-faktorisasi} mengapa selisih selama itu
   terjadi.
9. Sebuah kurva memiliki orde grup sekitar $2^{256}$. Berapa langkah yang
   dituntut algoritma rho Pollard untuk memecahkan ECDLP pada kurva itu?
   Nyatakan jawaban Anda dalam bentuk pangkat dua, lalu bandingkan dengan
   jumlah kunci AES-128 yang mungkin.
10. Bandingkan tiga kurva baku pada Tabel~\ref{tab:kurva-standar}. Mengapa
    Curve25519 tidak menuntut pemilihan titik basis seperti kurva NIST, dan
    mengapa hal itu justru dianggap keunggulan (Subbagian~\ref{subsec:efisiensi-ecc})?
\end{obereflection}

\section{Rangkuman}
\begin{itemize}
    \item Kurva eliptik $y^2 = x^3 + ax + b$ dengan $4a^3 + 27b^2 \neq 0$; himpunan titiknya bersama titik tak hingga $\mathcal{O}$ membentuk grup abelian $\langle G, + \rangle$.
    \item Penjumlahan titik $P + Q = R$ memakai gradien $m = (y_q - y_p)/(x_q - x_p)$, sedangkan penggandaan $2P$ memakai gradien garis singgung $m = (3x_p^2 + a)/(2y_p)$, dengan $x_r = m^2 - x_p - x_q$ dan $y_r = m(x_p - x_r) - y_p$.
    \item Pada medan berhingga $GF(p)$ semua operasi dilakukan modulo prima $p$, dan pembagian diganti perkalian dengan invers modulo, misalnya $m \equiv (3x_p^2 + a)(2y_p)^{-1} \bmod p$.
    \item Keamanan ECC bersandar pada \textbf{ECDLP}: diberikan $P$ dan $Q = kP$, sangat sukar menemukan skalar $k$ jika $k$ sangat besar.
    \item Kunci privat berupa integer $k \in [1, p-1]$ dan kunci publik berupa titik $Q = kP$; ECDH menghasilkan kunci bersama $K = abB$, sedangkan ECEG mengenkripsi pesan sebagai pasangan titik $[(kB), (P_M + kP_B)]$.
    \item ECC jauh lebih efisien daripada RSA: kunci 256 bit memberi keamanan setara RSA 3072 bit, sehingga cocok untuk \textit{smart card} dan perangkat IoT.
    \item Pesan teks harus dikodekan lebih dulu menjadi titik kurva, misalnya dengan metode Koblitz $x = mk + \text{offset}$ sampai diperoleh $y$ yang valid.
    \item Aritmetika kurva eliptik bertumpu pada aljabar abstrak. Himpunan
          titik pada kurva dengan aturan penjumlahan di atasnya memenuhi kelima
          aksioma grup (Tabel~\ref{tab:aksioma-grup-eliptik}), dan medan
          berhingga menyediakan invers bagi setiap unsur bukan nol
          (Tabel~\ref{tab:contoh-medan}).
    \item Penjumlahan $\Z$ bukan grup karena hanya $1$ dan $-1$ yang memiliki
          invers perkalian, sedangkan $\Z_n$ menjadi medan hanya bila $n$
          prima. Kenyataan itu menentukan modulus mana yang boleh dipakai.
    \item Di dalam $GF(2^m)$ penjumlahan dua polinom tidak lain adalah operasi
          XOR, sedangkan perkalian menuntut reduksi oleh polinom tak
          tereduksi berderajat $m$; sebagai contoh,
          $\code{57} \cdot \code{83} = \code{C1}$ di dalam $GF(2^8)$
          (Tabel~\ref{tab:gf2m-kali}).
    \item Titik di ketakhinggaan $\mathcal{O}$ memainkan dua peran sekaligus:
          elemen netral dan hasil penjumlahan dua titik berlawanan, sehingga
          $P + \mathcal{O} = P$ dan $P + (-P) = \mathcal{O}$
          (Persamaan~\eqref{eq:titik-o}).
    \item Rumus penjumlahan diturunkan dengan menyulihkan persamaan garis ke
          persamaan kurva. Hubungan Vieta memberi $x_p + x_q + x_r = m^2$,
          sehingga $x_r = m^2 - x_p - x_q$
          (Persamaan~\eqref{eq:koordinat-pq}); pada penggandaan akar $x_p$
          muncul dua kali dan $x_r = m^2 - 2x_p$
          (Persamaan~\eqref{eq:koordinat-penggandaan}).
    \item Bila ordinat sebuah titik nol, gradien garis singgungnya tidak
          terdefinisi dan $2P = \mathcal{O}$. Titik semacam itu berorde dua
          dan justru adalah akar-akar polinom kubik kurva.
    \item Perkalian skalar dihitung dengan \textit{double-and-add}, sehingga
          biaya menghitung $kP$ hanya sebanding dengan $\log_2 k$, bukan
          dengan $k$ (Persamaan~\eqref{eq:pelelaran}). Asimetri inilah yang
          membuat ECC praktis sekaligus aman.
    \item Banyaknya titik pada kurva tidak dapat ditentukan hanya dari nilai
          $p$; dalil Hasse hanya membatasinya pada selang
          $p + 1 \pm 2\sqrt{p}$. Karena itu orde grup dan orde titik basis
          harus diperiksa lebih dahulu, dan titik basis harus dipilih yang
          ordonya sama dengan orde grup.
    \item Enumerasi penuh memperlihatkan hal itu secara konkret. Kurva
          $y^2 \equiv x^3 + x + 6 \pmod{11}$ memiliki $12$ titik berhingga
          ditambah $\mathcal{O}$, sehingga orderya $13$
          (Tabel~\ref{tab:enumerasi-gf11} dan Tabel~\ref{tab:pelelaran-gf11}).
    \item Kunci publik ECC adalah sebuah titik, $Q = x \cdot B$
          (Persamaan~\eqref{eq:ecc-kunci-publik}). Dalam ECDH kedua pihak
          menghitung $K = a \cdot P_B = b \cdot P_A = abB$ tanpa pernah
          mengirim kunci privatnya (Persamaan~\eqref{eq:ecdh-kunci}).
    \item ECEG mengenkripsi plainteks $P_M$ menjadi pasangan titik
          $\bigl(kB,\ P_M + kP_B\bigr)$
          (Persamaan~\eqref{eq:eceg-cipher}). Dekripsinya berhasil karena
          $b(kB) = k(bB) = kP_B$, sehingga suku acak pada titik kedua
          terhapus (Persamaan~\eqref{eq:eceg-dekripsi}).
    \item Karena kunci sesaat $k$ dipilih baru untuk setiap pesan, enkripsi
          ECEG bersifat probabilistik: plainteks yang sama menghasilkan
          cipherteks yang berbeda pada setiap pengiriman. Sifat itu tidak
          dimiliki pemetaan karakter langsung, yang selalu memberi titik
          $P_M$ yang sama (Subbagian~\ref{subsec:pemetaan-langsung}).
    \item Metode Koblitz mengubah pesan menjadi titik dengan mencoba
          $x = mk + 1$, lalu $mk + 2$, dan seterusnya sampai ruas kanan kurva
          menjadi kuadrat modulo $p$ (Persamaan~\eqref{eq:koblitz-x});
          pendekodeannya memakai pembulatan ke bawah
          $m = \lfloor (x-1)/k \rfloor$
          (Persamaan~\eqref{eq:koblitz-dekode}).
    \item Algoritma \textit{index calculus} yang menyerang Diffie--Hellman
          dan ElGamal di atas bilangan bulat tidak dapat dipindahkan ke kurva
          eliptik. Serangan terbaik yang diketahui hanyalah rho Pollard, yang
          biayanya sebanding dengan akar orde grup
          (Tabel~\ref{tab:ecdlp-vs-faktorisasi}).
    \item Keunggulan panjang kunci itu hanya berlaku bagi kurva yang dipilih
          dengan benar. Kurva dengan orde yang habis dibagi bilangan kecil
          atau kurva yang dapat ditransfer ke medan berhingga biasa runtuh
          jauh lebih cepat, sehingga kurva baku seperti pada
          Tabel~\ref{tab:kurva-standar} harus dipakai bila memungkinkan.
\end{itemize}


\begin{obeassessment}
\textbf{Kuis Bab 14}
1. Uraikan aksioma grup yang harus dipenuhi oleh titik-titik pada kurva eliptik.
2. Bagaimana cara menghitung gradien $m$ untuk operasi penggandaan titik $2P$?
3. Mengapa ECC lebih cocok untuk perangkat mobile dibandingkan RSA?
\end{obeassessment}
\begin{competencychecklist}
\begin{tabularx}{\linewidth}{|X|c|}
\hline
Indikator Kompetensi & Tercapai \\
\hline
Saya dapat menjelaskan konsep grup dan medan dalam ECC & [ ] \\
\hline
Saya dapat menghitung penjumlahan dan penggandaan titik pada kurva & [ ] \\
\hline
Saya dapat menguraikan mekanisme ECDLP & [ ] \\
\hline
Saya dapat membedakan ECDH dan ECEG & [ ] \\
\hline
Saya dapat menganalisis trade-off panjang kunci RSA vs ECC & [ ] \\
\hline
\end{tabularx}
\end{competencychecklist}


\end{document}
